% ======================================================================
% Sections 7–8 — Limitations and Conclusion
% ======================================================================
\section{Limitations}\label{sec:limitations}

The study is a controlled evaluation on two public benchmarks, not a claim of
universal benefit. Several limitations bound the evidence.

\paragraph{Statistical power.}
All main results aggregate five training seeds. Sample standard deviations
and paired win counts are reported, and paired $t$-tests are given for
reference, but with $n=5$ the exact Wilcoxon signed-rank test cannot fall
below $p=0.0625$. We therefore avoid definitive significance claims and rely
on effect size, seed consistency, and the matched controls.

\paragraph{Scope of the controls.}
The matched-sparsity and capacity controls are restricted to Los-loop at
PH1. There is no matched-sparsity control on SZ-Taxi or at longer horizons,
no sparsified variant of the physical graph, and no sensitivity sweep over
DAGMA $\lambda_1$ or threshold parameters. The sparsity conclusion is
explicitly scoped to that single cell.

\paragraph{Datasets and horizons.}
Only two datasets are used, with different sampling intervals
($5$\,min and $15$\,min). Horizons are limited to $\mathrm{PH}\le 4$ at the
native sampling rate of each series. The $15$-minute Los-loop experiment
covers $15$--$60$ minutes ahead as a temporal-resolution variant; it is not
evidence for PH5--8 at $5$-minute sampling, and PH indices are not
comparable across sampling intervals.

\paragraph{Graph-learning assumptions.}
Graphs are estimated with linear DAGMA ($\ell_2$ score, acyclicity
regularizer). Nonlinear dependency, measurement noise beyond the linear SEM
assumption, and misspecification of the lag depth $L=3$ are not examined.
Measured runtimes are non-negligible for the multi-lag fit (on the order of
hours for the $828$-variable construction on Los-loop), which limits rapid
replication at larger $N$.

\paragraph{Static graphs; gated use.}
The learned graphs are fitted once and remain static. The Mix mechanism
varies how those static graphs are consumed per node and timestep; it does
not produce time-varying edge sets. Genuinely time-varying learned graphs
remain future work.

\paragraph{Backbone scope.}
Experiments cover the GCN and T-GCN family used here. Other spatio-temporal
architectures (attention-based models, diffusion variants, large
transformers) are outside the study.

\section{Conclusion}\label{sec:conclusion}

We asked when and how learned graph structure helps graph-based traffic
forecasting, separating the \emph{source} of the graph from the
\emph{mechanism that consumes it}. Across two public benchmarks, twelve
configurations, four horizons, and five seeds, three conclusions follow.
First, the dense physical road-network graph is not a sound default: a
graph-free identity control outperforms it in every dataset--horizon cell,
and single learned contemporaneous graphs improve on the physical graph but
never on the graph-free control. Second, lag-specific multi-lag DAGMA
graphs are useful on Los-loop only when consumed in a temporally aligned
manner; feeding the identical graphs as a static union is markedly worse,
so consumption is at least as important as learning, while a global lag
mixture is insufficient and a per-node, per-timestep gate adds a further
seed-consistent gain (up to $14.5\%$ relative RMSE reduction versus the
graph-free baseline at PH1). Third, the benefit is conditional: on SZ-Taxi,
where the learned multi-lag graphs retain almost no edges, all learned
variants remain within seed variability of the graph-free baseline, and
matched-budget controls show that sparsity alone does not explain the
Los-loop gain.

Future work includes longer horizons at $5$-minute resolution, additional
networks, sparsified-physical controls, $\lambda$/threshold sensitivity,
other backbones, and genuinely time-varying graph structure. Practically,
we recommend that graph-based traffic forecasters report a graph-free
control and state explicitly how any learned graph is consumed over time.
